\documentclass[12pt]{article}

\usepackage[utf8]{inputenc}
\usepackage[spanish]{babel}
\usepackage[shortlabels]{enumitem}

\usepackage{mathtools}
\usepackage{amssymb}
\usepackage{amsthm}

\usepackage{geometry}
\geometry{margin=2.5cm}

\usepackage{hyperref}

\title{Ejercicios Aplicaciones Lineales}
\author{}

\begin{document}
\maketitle

\section*{Ejercicio 1}

Dado el endomorfismo de \(\mathbb{R}^3\) definido por las ecuaciones
\[
(y_1,y_2,y_3)=(x_1+x_2+x_3,\;x_1+x_2-x_3,\;x_3),
\]
obtener:

\begin{enumerate}
    \item[a)] Núcleo e imagen de la aplicación lineal.
    
    \item[b)] Clasificar la aplicación lineal.
    
    \item[c)] \(f(V)\), siendo
    \[
    V=\{(x_1,x_2,x_3)\in\mathbb{R}^3:x_1+x_2+x_3=0\}.
    \]
\end{enumerate}

\section*{Ejercicio 2}

Calcular núcleo e imagen de la aplicación lineal
\[
f:\mathbb{R}^4\longrightarrow \mathbb{R}^3
\]
definida por
\[
f(x_1,x_2,x_3,x_4)=(x_1+x_2,\;0,\;x_1-x_2).
\]

\section*{Ejercicio 3}

Consideramos el endomorfismo
\[
f:P_2(x)\longrightarrow P_2(x)
\]
definido sobre el espacio vectorial de polinomios de grado menor o igual que \(2\) con coeficientes reales, que verifica:
\[
\ker f=\{a+bx+cx^2\in P_2(x): a=b=c\}
\]
y tal que los vectores que tienen término independiente nulo se transforman en sí mismos.

\begin{enumerate}
    \item[a)] Determinar la matriz \(M_{B_C}(f)\) del endomorfismo en la base canónica.
    
    \item[b)] Obtener la expresión implícita del subespacio \(\operatorname{Im} f\) y razonar de dos formas distintas cuál es su dimensión.
    
    \item[c)] Si
    \[
    S=\{(\lambda-2\mu+\delta,\lambda-2\mu-\delta,\lambda-2\mu): \lambda,\mu,\delta\in\mathbb{R}\},
    \]
    determinar en forma implícita \(f(S)\) y razonar si tiene sentido la dimensión obtenida.
    
    \item[d)] Dada la base de \(P_2(x)\)
    \[
    B^*=\{1+x+x^2,\;1+x,\;1\},
    \]
    obtener por coordenadas y matricialmente la matriz \(M_{B^*}(f)\).
\end{enumerate}

\section*{Ejercicio 4}

Dadas las aplicaciones lineales
\[
f:\mathbb{R}^2\longrightarrow \mathbb{R}^3
\]
tal que
\[
f(1,0)=(2,1,2),
\qquad
f(0,1)=(3,2,1),
\]
y
\[
g:\mathbb{R}^3\longrightarrow \mathbb{R}^2
\]
tal que
\[
g(1,0,0)=(3,-1),
\qquad
g(0,1,0)=(1,0),
\qquad
g(0,0,1)=(1,-2).
\]

Si definimos
\[
h=g\circ f,
\]
calcular \(h(4,-2)\) matricialmente y por coordenadas.

\section*{Ejercicio 5}

Dada la aplicación lineal
\[
f:\mathbb{R}^3\longrightarrow \mathbb{R}^2
\]
tal que
\[
f(x,y,z)=(2x-y,2y+z),
\]
y la aplicación lineal
\[
g:\mathbb{R}^2\longrightarrow \mathbb{R}^3
\]
tal que
\[
g(x,y)=(4x+2y,y,x+y).
\]

Consideramos en \(\mathbb{R}^3\) la base
\[
B=\{(1,0,0),(1,1,0),(1,1,1)\},
\]
y en \(\mathbb{R}^2\) la base
\[
B'=\{(1,0),(1,1)\}.
\]

\begin{enumerate}
    \item[a)] Determinar \(M_{CC'}(f)\) y \(M_{C'C}(g)\), siendo \(C\) y \(C'\) respectivamente las bases canónicas de \(\mathbb{R}^3\) y \(\mathbb{R}^2\).

    \item[b)] Obtener una base de \(\ker f\) y otra base de \(\operatorname{Im} g\). Clasificar las aplicaciones \(f\) y \(g\).

    \item[c)] Calcular
    \[
    [f(2,0,-1)]_{B'}
    \]
    utilizando la matriz \(M_{BB'}(f)\) y comprobar el resultado en las bases canónicas.
\end{enumerate}

\section*{Ejercicio 6}

Consideramos la aplicación lineal
\[
f:\mathbb{R}^3\longrightarrow \mathbb{R}^2
\]
de la que sabemos que:
\[
\ker f=\{(x,y,z)\in\mathbb{R}^3:x+ay+z=0,\ ax+y+z=0\},
\]
y además:
\[
f(0,0,a-1)=(0,a-1),
\qquad
f(1,0,1)=(2,1),
\]
donde \(a\in\mathbb{R}\).

\begin{enumerate}
    \item[a)] Obtener dimensiones y bases de \(\ker f\) y de \(\operatorname{Im} f\) según los valores de \(a\).

    \item[b)] Calcular la matriz de la aplicación lineal en las bases canónicas de \(\mathbb{R}^3\) y \(\mathbb{R}^2\).

    \item[c)] Para \(a=1\), hallar, si existen, los vectores que se transforman en el vector \((2,1)\) y en el vector \((2,3)\).
\end{enumerate}

\end{document}